\documentclass[a4paper,12pt]{article}

\usepackage[utf8]{inputenc}
\usepackage[ukrainian]{babel}
\usepackage{amsmath, amssymb}

\title{Задача про ефективне розрізнення центра та фокуса}
\author{}
\date{}

\begin{document}

\maketitle

\section*{Постановка задачі}

Розглянемо автономну систему диференціальних рівнянь на площині:
\begin{equation}
\begin{cases}
\dot{x} = y + q(x,y), \\
\dot{y} = -x + p(x,y),
\end{cases}
\end{equation}
де $p(x,y)$ та $q(x,y)$ --- аналітичні функції, що мають розклади у степеневі ряди, починаючи з членів другого порядку:
\[
p(x,y) = \sum_{k=2}^{\infty} p_k(x,y), \quad
q(x,y) = \sum_{k=2}^{\infty} q_k(x,y),
\]
де $p_k(x,y), q_k(x,y)$ --- однорідні многочлени степеня $k$.

Нульова точка $(0,0)$ є положенням рівноваги.

\section*{Проблема центра--фокуса}

Необхідно визначити, чи є точка $(0,0)$:
\begin{itemize}
\item \textbf{центром} (усі достатньо близькі траєкторії замкнуті),
\item чи \textbf{фокусом} (траєкторії мають спіральний характер).
\end{itemize}

Відомо, що це еквівалентно виконанню нескінченної системи умов:
\[
V_1 = 0, \quad V_2 = 0, \quad V_3 = 0, \quad \dots
\]
де $V_k$ --- так звані \textit{фокусні величини} (або коефіцієнти Ляпунова), які обчислюються через коефіцієнти $p_k, q_k$.

\section*{Основна задача (ефективне розрізнення)}

Нехай праві частини системи є многочленами степеня $n$.

Потрібно:

\begin{quote}
Знайти (або довести існування) такого числа $N(n)$, що якщо виконуються умови
\[
V_1 = V_2 = \dots = V_{N(n)} = 0,
\]
то автоматично виконуються всі наступні умови:
\[
V_k = 0 \quad \text{для всіх } k > N(n).
\]
\end{quote}

\section*{Мета}

\begin{itemize}
\item Побудувати оцінку або явний вигляд $N(n)$.
\item Або довести, що такого універсального числа $N(n)$ не існує.
\end{itemize}

\section*{Зауваження}

\begin{itemize}
\item Для окремих класів систем (наприклад, квадратичних, $n=2$) такі умови відомі.
\item У загальному випадку задача залишається відкритою.
\end{itemize}

\end{document}